\section{Front-tracking em tr\^es dimens\~oes}
\label{sec:tridim}

\subsection{O que sobe de dimens\~ao, e o que n\~ao}

A maquinaria de \emph{transfer\^encia} --- a parte dif\'icil --- j\'a era
tridimensional e verificada. O m\'odulo da fronteira imersa trata superf\'icie em
3D como \textbf{codimens\~ao 1} com peso igual a \'area, e a lei do espalhamento
j\'a estava escrita como $dV = \text{peso}\cdot h^{DIM-1}$, gen\'erica. Busca de
faceta na malha escalonada, n\'ucleo de Roma e interpola\c{c}\~ao vieram de l\'a
sem uma linha de altera\c{c}\~ao, e o \texttt{example3d\_Schaefer\-Turek} j\'a os
exercitava.

O que n\~ao sobe \'e o n\'ucleo da frente, e n\~ao por descuido:

\begin{adjustbox}{max width=\linewidth}
\begin{tabular}{>{\raggedright\arraybackslash}p{3.2cm}>{\raggedright\arraybackslash}p{4.6cm}>{\raggedright\arraybackslash}p{6.0cm}}
\toprule
& 2D hoje & 3D exige \\
\midrule
conectividade & impl\'icita no \'indice, $i \to i+1$ & malha de tri\^angulos expl\'icita, com adjac\^encia por aresta \\
medida fechada & f\'ormula do la\c{c}o & soma de tri\^angulos, e \textbf{volume} pelo divergente \\
curvatura & circunc\'irculo de tr\^es pontos & n\~ao generaliza \\
cirurgia & inserir e remover ponto & dividir, colapsar e \textbf{girar} aresta preservando a variedade \\
\bottomrule
\end{tabular}
\end{adjustbox}

Da\'i a decis\~ao de m\'odulo separado: os dois n\'ucleos quase n\~ao t\^em
c\'odigo em comum, e o bidimensional est\'a verificado fechando um
\emph{benchmark}. Os dois arquivos se recusam a compilar na dimens\~ao do outro.

\subsection{Tens\~ao superficial sem curvatura}

A for\c{c}a sai da integral de linha na borda de cada tri\^angulo,
\begin{equation}
  \mathbf{F} = \sigma \oint (\mathbf{t}\times\mathbf{n})\, ds ,
  \label{eq:tensao3d}
\end{equation}
que para aresta reta de $a$ a $b$ num tri\^angulo de normal $\mathbf{n}$ vale
$(b-a)\times\mathbf{n}$. Curvatura nunca \'e calculada: ela entra pela
\emph{diferen\c{c}a de normais} entre os dois tri\^angulos que compartilham cada
aresta interior, sem nunca ser nomeada.

\paragraph{DOIS sobre $R$, e n\~ao um.} Em 2D o salto de Laplace \'e $\sigma/R$;
em 3D a curvatura m\'edia da esfera \'e $2/R$ e o salto \'e $2\sigma/R$. A
Se\c{c}\~ao~\ref{sec:balanco} fechou o caso bidimensional a $0{,}02\%$ contra
$\sigma/R$, e reaproveitar aquele n\'umero aqui validaria o errado. A
advert\^encia est\'a no cabe\c{c}alho do m\'odulo, no arreio e no exemplo ---
tr\^es lugares, porque a tenta\c{c}\~ao \'e real.

\subsection{Os degraus, e o or\'aculo de cada um}

Nenhum degrau foi dado sem or\'aculo fechado, e nenhum c\'odigo entrou sem o
teste que o exercita:

\begin{adjustbox}{max width=\linewidth}
\pgfplotstabletypeset[
  string type,
  columns/degrau/.style={column name={degrau}},
  columns/oraculo/.style={column name={or\'aculo}},
  columns/resultado/.style={column name={melhor medido}},
  columns/ordem/.style={column name={ordem}},
  every head row/.style={before row=\toprule, after row=\midrule},
  every last row/.style={after row=\bottomrule},
]{\dat{tridim-degraus.dat}}
\end{adjustbox}

\paragraph{A ordem do B2 \'e um, e isso n\~ao \'e defeito.} Dobrar a
resolu\c{c}\~ao (com $ds/h$ constante) derruba o erro do salto de Laplace por
$1{,}74$, e a corrente esp\'uria por $1{,}59$. \'E primeira ordem, e \'e o
esperado: o delta regularizado borra a interface sobre $\sim 3h$, e m\'etodo de
fronteira imersa com esse ingrediente \'e de primeira ordem no salto.
\textbf{Dois pontos n\~ao estabelecem ordem}, e o terceiro custaria $512$ mil
c\'elulas --- est\'a dito, e n\~ao foi convertido em expoente.

\subsection{A fra\c{c}\~ao de volume, e por que n\~ao \'e recorte exato}

O caminho bidimensional recorta o pol\'igono contra a caixa
(Sutherland--Hodgman) e aplica a f\'ormula do la\c{c}o. Isso \textbf{n\~ao
generaliza}: recortar superf\'icie triangulada n\~ao-convexa contra caixa \'e
outro problema, e fechar os la\c{c}os nas faces da caixa \'e onde defeito mora.

O que se faz \'e distinto e mais barato: longe da interface a caixa est\'a
inteira de um lado, decidido por dist\^ancia contra a meia-diagonal; cortada, a
interface vira o \emph{plano} que passa na dist\^ancia medida com a normal
medida, e o volume sai da f\'ormula fechada de plano--caixa. Exato para
interface plana, $O(h^2)$ para suave --- a mesma aproxima\c{c}\~ao que o PLIC faz
no VOF.

O sinal \textbf{n\~ao} vem da normal: a normal do tri\^angulo mais pr\'oximo erra
o lado perto de aresta e v\'ertice, e o sinal sai de \emph{paridade de
cruzamentos} num raio. Dist\^ancia e sinal s\~ao medidos por caminhos
independentes de prop\'osito.

\paragraph{O or\'aculo que vale mais que a soma.} Somar $f\,h^3$ sobre uma
parti\c{c}\~ao e comparar com o volume fechado \'e necess\'ario mas fraco: erros
que se cancelam entre c\'elulas passam. O arreio compara tamb\'em, \textbf{c\'elula
a c\'elula}, contra uma subamostragem de mil pontos que n\~ao usa plano nenhum
--- pior diferen\c{c}a $0{,}029$. Concordar com um m\'etodo de outra natureza \'e
evid\^encia de outra esp\'ecie.

\subsection{A grade espacial: acelerar sem aproximar}

A consulta \'e $O(n_t)$ por c\'elula, e preencher a fra\c{c}\~ao custava
$0{,}09$~s por passo com $8$ mil c\'elulas contra $1280$ tri\^angulos. Uma grade
uniforme de tri\^angulos resolve, com um contrato explicito: ela restringe
\emph{quais} tri\^angulos s\~ao testados, \textbf{nunca o que se calcula}.

O arreio compara os dois caminhos em $20$ mil consultas de dentro/fora, $2$ mil
de dist\^ancia e nas $8$ mil c\'elulas de uma parti\c{c}\~ao, e exige que n\~ao
difiram \textbf{em bit nenhum}. \'E a \'unica forma de uma acelera\c{c}\~ao n\~ao
virar uma aproxima\c{c}\~ao silenciosa. Resultado: zero diferen\c{c}as, $14{,}1\times$
no arreio e $37\times$ no solver.

\paragraph{E o gargalo estava onde eu n\~ao olhava.} O ganho parou em
$2{,}2\times$ at\'e eu descobrir que \texttt{\_fracao} recalculava a
\emph{envolt\'oria} a cada c\'elula, varrendo os $642$ v\'ertices --- cinco
milh\~oes de opera\c{c}\~oes que a grade n\~ao elimina, porque ela acelera busca
por \emph{tri\^angulo} e n\~ao varredura de \emph{v\'ertice}. Quem apontou o
lugar certo foi uma medida \emph{negativa}: o atalho de ocupa\c{c}\~ao, que eu
esperava que resolvesse, n\~ao mudou nada.

\subsection{Adve\c{c}\~ao acoplada}

A superf\'icie anda com a velocidade interpolada da malha. A ordem dentro do
passo \'e a mesma do caminho bidimensional: move com $\mathbf{u}^n$ --- j\'a
projetada, e portanto discretamente livre de diverg\^encia ---, faz a cirurgia, e
espalha a for\c{c}a nas posi\c{c}\~oes \emph{novas}.

Dois or\'aculos. O \textbf{volume}, porque campo livre de diverg\^encia preserva
o volume fechado; e a \textbf{parti\c{c}\~ao da unidade}, porque os pesos de Roma
somam $1$ em aritm\'etica exata e soma diferente de $1$ significa suporte
incompleto. Nesse caso o adaptador \textbf{aborta} em vez de devolver zero ---
devolver zero pregaria o v\'ertice no lugar em sil\^encio, que \'e o defeito que
a Se\c{c}\~ao~\ref{sec:posse} levou meio dia para achar em 2D.

\paragraph{E o guarda disparou.} Num caso de $Ca=20$ a gota foi carregada pelo
v\'ortice at\'e $y=0{,}925$, e o suporte do n\'ucleo ($\pm 1{,}5h$) alcan\c{c}ou
a fronteira do dom\'inio. O guarda abortou. N\~ao \'e defeito da adve\c{c}\~ao:
\'e li\c{c}\~ao de \emph{montagem de caso} --- v\'ortice de cavidade transporta
para a tampa, e trocar a gota de $0{,}25$ para $0{,}15$ n\~ao mudou nada, porque
o tamanho n\~ao era a vari\'avel.

At\'e o passo em que abortou, a combina\c{c}\~ao adve\c{c}\~ao mais cirurgia
ficou exercitada: $130$ opera\c{c}\~oes de cirurgia, $\chi=2$ o tempo todo,
parti\c{c}\~ao da unidade em $6\times10^{-15}$, e deriva de volume limitada a
$5{,}8\times10^{-4}$ --- que \emph{sobe e volta}, em vez de acumular.

\subsection{Sete defeitos, e o or\'aculo que pegou cada um}

Nenhum destes foi achado lendo o c\'odigo:

\begin{adjustbox}{max width=\linewidth}
\pgfplotstabletypeset[
  string type,
  columns/defeito/.style={column name={defeito}},
  columns/comooraculopegou/.style={column name={como apareceu}},
  every head row/.style={before row=\toprule, after row=\midrule},
  every last row/.style={after row=\bottomrule},
]{\dat{tridim-defeitos.dat}}
\end{adjustbox}

Tr\^es merecem nota. O \textbf{sinal da for\c{c}a} n\~ao foi trocado at\'e o
or\'aculo passar: comparei contra o gradiente \emph{num\'erico} da \'area e o
c\'odigo dava exatamente $+\sigma\nabla A$; a for\c{c}a \'e menos o gradiente,
porque tens\~ao minimiza \'area. A raz\~ao de estar certo \'e a variacional, e
n\~ao o or\'aculo.

O \textbf{giro de aresta} foi atribu\'ido por medida e n\~ao por suspeita: com um
seletor de etapas, divis\~ao sozinha dava zero invertidos, colapso zero,
divis\~ao mais colapso zero, e \emph{qualquer} combina\c{c}\~ao contendo o giro
dava de $1$ a $15$.

E o \textbf{ponto m\'edio da corda}: dividir aresta inseria no meio da corda, que
numa superf\'icie convexa fica para dentro, e $1500$ divis\~oes custavam
$5{,}2\%$ do volume. O que isolou a causa foi refinar o passo de tempo pela
metade e o erro \emph{n\~ao mudar} --- prova de que n\~ao era o integrador. O
conserto \'e de \emph{coloca\c{c}\~ao}, e n\~ao corre\c{c}\~ao de volume imposta:
o ponto novo vai no meio da curva de Hermite c\'ubica que interpola os extremos
com as normais. For\c{c}ar a conserva\c{c}\~ao teria tornado o or\'aculo vazio.

\subsection{O que est\'a estabelecido, e o que n\~ao}

Est\'a estabelecido: geometria, topologia e for\c{c}a contra a esfera anal\'itica;
adve\c{c}\~ao cinem\'atica revers\'ivel com ordem medida; salto de Laplace a
$0{,}19\%$; fra\c{c}\~ao de volume em segunda ordem, cruzada com m\'etodo
independente; grade acelerando $37\times$ sem mudar um bit; e adve\c{c}\~ao
acoplada com cirurgia, preservando topologia.

\subsection{A deriva \'e discretiza\c{c}\~ao, e a cirurgia a reduz}
\label{sec:tridim-duasres}

A pergunta que esta se\c{c}\~ao carregava em aberto --- se a deriva de volume da
adve\c{c}\~ao acoplada tem componente sistem\'atica --- foi fechada por dois
experimentos controlados, cada um variando \textbf{uma} coisa.

\paragraph{Duas resolu\c{c}\~oes.} Mesmo caso ($R=0{,}25$, $\sigma=0{,}2$, tampa
$4{,}0$, $\Delta t = 5\times10^{-4}$, at\'e $t=0{,}25$) em duas malhas, com a
resolu\c{c}\~ao da superf\'icie acompanhando ($ds/h = 0{,}75$ nos dois pares:
$N=20$ com $nsub=3$, $N=40$ com $nsub=4$). Sem isso seriam duas vari\'aveis, e o
erro de deixar $nsub$ fixo j\'a custou uma s\'erie inteira na
Se\c{c}\~ao~\ref{sec:conv-serieC}.

\begin{adjustbox}{max width=\linewidth}
\pgfplotstabletypeset[
  string type,
  columns/passo/.style={column name={passo}},
  columns/aa0/.style={column name={$A/A_0$}},
  columns/n20/.style={column name={$N=20$}},
  columns/cir20/.style={column name={cir.}},
  columns/n40/.style={column name={$N=40$}},
  columns/cir40/.style={column name={cir.}},
  columns/razao/.style={column name={raz\~ao}},
  every head row/.style={before row=\toprule, after row=\midrule},
  every last row/.style={after row=\bottomrule},
]{\dat{tridim-duasresolucoes.dat}}
\end{adjustbox}

A raz\~ao fica entre $3{,}2$ e $3{,}7$ ao longo de \emph{toda} a corrida ---
segunda ordem daria $4$, primeira daria $2$. Nos tr\^es pontos anteriores \`a
primeira cirurgia, que s\~ao adve\c{c}\~ao pura contra adve\c{c}\~ao pura, a
raz\~ao \'e $3{,}63$, $3{,}70$ e $3{,}61$: m\'edia $3{,}65$, ou ordem observada
$\mathbf{1{,}87}$.

\textbf{A deriva \'e erro de discretiza\c{c}\~ao e some com refino.} N\~ao h\'a
componente sistem\'atica: vi\'es de interpola\c{c}\~ao n\~ao cairia com $h$.

Uma ressalva que a pr\'opria tabela exibe: no passo $500$ a raz\~ao cai a
$3{,}21$, e ali o caso fino j\'a fez $252$ opera\c{c}\~oes de cirurgia contra
$55$ do grosso --- a compara\c{c}\~ao deixa de ser s\'o de adve\c{c}\~ao. Os
pontos limpos s\~ao os primeiros, e s\~ao eles que pesam.

\paragraph{Com e sem cirurgia.} O segundo experimento desliga a cirurgia e deixa
todo o resto id\^entico:

\begin{adjustbox}{max width=\linewidth}
\pgfplotstabletypeset[
  string type,
  columns/passo/.style={column name={passo}},
  columns/comcirurgia/.style={column name={com cirurgia}},
  columns/semcirurgia/.style={column name={sem cirurgia}},
  columns/razao/.style={column name={pior por}},
  every head row/.style={before row=\toprule, after row=\midrule},
  every last row/.style={after row=\bottomrule},
]{\dat{tridim-cirurgia.dat}}
\end{adjustbox}

\textbf{A cirurgia n\~ao causa a deriva: ela a reduz em $42\%$.} Era o suspeito
nomeado --- a coloca\c{c}\~ao do ponto novo j\'a havia custado $5{,}2\%$ de
volume no B1 antes do conserto de Hermite --- e o experimento o inocenta. Sem
cirurgia os tri\^angulos esticam e a superf\'icie deformada fica pior
representada.

As duas corridas s\~ao \textbf{id\^enticas at\'e o passo $200$}, que \'e onde a
primeira cirurgia acontece. \'E a verifica\c{c}\~ao de que a \'unica diferen\c{c}a
entre elas \'e a que se quis introduzir.

\subsection{O que continua em aberto}

O \textbf{paralelo} em 3D.

O espalhamento e a
interpola\c{c}\~ao j\'a est\~ao escritos na forma em lote, com o filtro de posse,
porque o caminho bidimensional aprendeu isso depois e pagou a reescrita --- mas
em $n_p>1$ nada foi medido. E o bif\'asico com salto de propriedade: o
encanamento est\'a verificado com $\rho$ e $\mu$ \emph{iguais} nas duas fases, que
\'e o est\'agio que move o encanamento sem mudar a f\'isica.
